\documentclass[11pt]{article}
\usepackage[margin=1in]{geometry}
\usepackage{amsmath,amssymb,amsthm}
\usepackage{graphicx}
\usepackage{booktabs}
\usepackage{microtype}
\usepackage[hidelinks]{hyperref}
\usepackage{xcolor}

\newtheorem{proposition}{Proposition}
\newtheorem{remark}{Remark}

% ---- result macros (filled by scripts/make_figures.py output) ----
\newcommand{\capmapegbr}{0.31\%}
\newcommand{\capmapenn}{1.75\%}
\newcommand{\gerrgbr}{0.70\%}
\newcommand{\gerrnn}{4.03\%}
\newcommand{\chierrgbr}{1.71\%}
\newcommand{\chierrnn}{28.7\%}

\title{\textbf{Ansatz: Cost-Model Routing for Verified Electromagnetic
Simulation in Superconducting Qubit Design}}
\author{Ansatz AI\\ \texttt{contact@useansatz.ai}}
\date{August 2026 \quad (v0 technical report)}

\begin{document}
\maketitle

\begin{abstract}
The inner loop of superconducting qubit design is electromagnetic simulation:
each geometry iterate requires a multi-conductor capacitance extraction whose
lumped-oscillator reduction yields the Hamiltonian parameters designers
actually target. We describe the v0 Ansatz engine, which accelerates this loop
with two coupled components. First, a forward surrogate trained on 1{,}934
experimentally validated Ansys Q3D capacitance matrices (the SQuADDS database)
predicts capacitance matrices at microsecond latency, roughly $5\times$ more
accurately than nearest-design lookup in distribution ($\capmapegbr$ vs.\
$\capmapenn$ mean absolute percentage error), with downstream coupling and
dispersive-shift errors of $\gerrgbr$ and $\chierrgbr$. Second, because every
data-driven predictor degrades outside its training hull, a \emph{routed,
verified} field-solver answers arbitrary designs: a cost-model router chooses
per instance among sparse-direct, AMG-preconditioned Krylov,
surrogate-initialized matrix-free Krylov, and hybrid pipelines, every pipeline
terminates in residual verification, and a runtime escalation monitor bounds
the cost of routing mistakes. On geometry distributions anchored to the same
validated design family, the routed solver tracks the per-instance oracle
within measurement noise and dominates every fixed pipeline choice as
resolution varies, while HINTS-style fixed hybrid schedules are never
competitive. We argue this instance-level, verification-first reformulation of
learned solver routing --- rather than per-iteration classification or
uncertainty-quantified surrogation alone --- is the form in which neural
acceleration becomes deployable in engineering design loops.
\end{abstract}

\section{Introduction}
\label{sec:intro}
% Design loop context; cost of Q3D/HFSS in the loop; why surrogates alone are
% not deployable (silent extrapolation failure); why verification-first.
TODO

\section{Setting}
\label{sec:setting}
% Planar transmon family (xmon + claw); electrostatics; capacitance matrix;
% LOM reduction to (E_C, f_q, alpha, g, chi). Real data: SQuADDS.
TODO

\section{Method}
\label{sec:method}
\subsection{Tier 1: forward surrogate on validated Q3D data}
TODO
\subsection{Tier 2: routed, verified field solves}
% pipelines, shared per-design context, cost-model router, escalation monitor,
% unconditional verification. Contrast: per-iteration greedy routing
% (arXiv:2509.24814), HINTS fixed schedules.
TODO

\subsection{Guardrail properties}
\begin{proposition}[Verification safety]
Every field returned by any pipeline satisfies
$\|b - Au\|_2 / \|b\|_2 \le \tau$ on the target discretization; on
verification failure the engine falls back to the sparse direct solve, so the
routed result equals the classical result in the worst case.
\end{proposition}

\begin{proposition}[Bounded escalation regret]
Let $t_s$ be the time spent in a monitored pipeline before escalation fires,
and $T_{\mathrm{alt}}$ the (warm-started) time of the escalation target. The
routed cost is at most $t_s + T_{\mathrm{alt}}$, where $t_s$ is bounded by the
monitor's projection rule; in particular a stalled pipeline cannot consume
more than the patience window beyond its projected-competitive segment.
\end{proposition}

\section{Experiments}
\label{sec:experiments}
\subsection{Forward model on real Q3D data}
% table: capacitance MAPE + downstream errors, iid and out-of-hull
TODO
\subsection{Routed solver benchmarks}
% per-resolution pipeline times; router vs oracle vs fixed; OOD; choices
TODO
\subsection{Ablations}
% escalation on/off; surrogate quality vs routing benefit; decision overhead
TODO

\section{Related work}
% HINTS; greedy PDE router; SQuADDS interpolation; AMG practice; neural
% operators for EM; verification-first hybrids.
TODO

\section{Limitations and roadmap}
% 2D plan-view v0; calibration to 3D; Palace backend; eigenmode analyses;
% cost-model transfer across hardware.
TODO

\bibliographystyle{plain}
\begin{thebibliography}{9}
\bibitem{squadds} S.~Shanto et al. SQuADDS: A validated design database and
simulation workflow for superconducting qubit design. \emph{Quantum} 8, 1465
(2024).
\bibitem{hints} E.~Zhang et al. Blending neural operators and relaxation
methods in PDE numerical solvers. \emph{Nature Machine Intelligence} (2024).
\bibitem{router} S.~Rayan, Y.~Patel, A.~Tewari. A greedy PDE router for
blending neural operators and classical methods. arXiv:2509.24814 (2025).
\bibitem{koch} J.~Koch et al. Charge-insensitive qubit design derived from the
Cooper pair box. \emph{Phys. Rev. A} 76, 042319 (2007).
\bibitem{pyamg} N.~Bell et al. PyAMG: Algebraic multigrid solvers in Python.
\emph{JOSS} 8(87), 5495 (2023).
\bibitem{minev} Z.~Minev et al. Circuit quantum electrodynamics (cQED) with
modular quasi-lumped models. arXiv:2103.10344 (2021).
\end{thebibliography}

\end{document}
